\chapter{2024年北京卷}

\section{单选题}

\begin{problem}
已知集合 \(M=\{x\mid -3<x<1\}\)，\(N=\{x\mid -1\le x<4\}\)，则 \(M\cup N=\)（\quad）

\choices
    {\(\{x\mid -1\le x<1\}\)}
    {\(\{x\mid x>-3\}\)}
    {\(\{x\mid -3<x<4\}\)}
    {\(\{x\mid x<4\}\)}

\end{problem}

\begin{answer}
C
\end{answer}

\begin{solution}
根据题意可得
\(M\cup N=\{x\mid -3<x<4\}\).
故选 C.
\end{solution}
\begin{problem}
已知 \(\frac{z}{\i}=-1-\i\)，则 \(z=\)（\quad）

\choices
    {\(-1-\i\)}
    {\(-1+\i\)}
    {\(1-\i\)}
    {\(1+\i\)}

\end{problem}

\begin{answer}
C
\end{answer}

\begin{solution}
由 \(\frac{z}{\i}=-1-\i\) 得
\(z=\i\left( -1-\i \right)=1-\i\).
故选 C.
\end{solution}
\begin{problem}
圆 \(x^2+y^2-2x+6y=0\) 的圆心到直线 \(x-y+2=0\) 的距离为（\quad）

\choices
    {\(\sqrt{2}\)}
    {2}
    {3}
    {\(3\sqrt{2}\)}

\end{problem}

\begin{answer}
D
\end{answer}

\begin{solution}
将圆的方程配方得
\(\left( x-1 \right)^2+\left( y+3 \right)^2=10\)，
所以圆心为 \((1,-3)\). 于是圆心到直线 \(x-y+2=0\) 的距离为
\(\frac{|1-(-3)+2|}{\sqrt{1^2+(-1)^2}}=\frac{6}{\sqrt{2}}=3\sqrt{2}\).
故选 D.
\end{solution}
\begin{problem}
在 \(\left( x-\sqrt{x} \right)^4\) 的展开式中，\(x^3\) 的系数为（\quad）

\choices
    {6}
    {\(-6\)}
    {12}
    {\(-12\)}

\end{problem}

\begin{answer}
A
\end{answer}

\begin{solution}
\(\left( x-\sqrt{x} \right)^4\) 的二项展开式的通项为
\(T_{r+1}=\mathrm{C}_4^r x^{4-r}\left( -\sqrt{x} \right)^r=\mathrm{C}_4^r (-1)^r x^{4-\frac{r}{2}}\)（\(r=0,1,2,3,4\)）.
令 \(4-\frac{r}{2}=3\)，得 \(r=2\). 因此 \(x^3\) 的系数为
\(\mathrm{C}_4^2(-1)^2=6\).
故选 A.
\end{solution}
\begin{problem}
设 \(\bs{a},\bs{b}\) 是向量，则“\(\left( \bs{a}+\bs{b} \right)\cdot\left( \bs{a}-\bs{b} \right)=0\)”是“\(\bs{a}=-\bs{b}\) 或 \(\bs{a}=\bs{b}\)”的（\quad）

\choices
    {充分不必要条件}
    {必要不充分条件}
    {充要条件}
    {既不充分也不必要条件}

\end{problem}

\begin{answer}
B
\end{answer}

\begin{solution}
因为
\(\left( \bs{a}+\bs{b} \right)\cdot\left( \bs{a}-\bs{b} \right)=|\bs{a}|^2-|\bs{b}|^2\)，
所以
\(\left( \bs{a}+\bs{b} \right)\cdot\left( \bs{a}-\bs{b} \right)=0\)
等价于 \(|\bs{a}|=|\bs{b}|\). 若 \(\bs{a}=\bs{b}\) 或 \(\bs{a}=-\bs{b}\)，则必有 \(|\bs{a}|=|\bs{b}|\)，所以必要性成立.

反过来，只由 \(|\bs{a}|=|\bs{b}|\) 不能推出 \(\bs{a}=\bs{b}\) 或 \(\bs{a}=-\bs{b}\). 例如 \(\bs{a}=(1,0)\)，\(\bs{b}=(0,1)\) 时，\(|\bs{a}|=|\bs{b}|\)，但 \(\bs{a}\ne \bs{b}\)，且 \(\bs{a}\ne -\bs{b}\).

故选 B.
\end{solution}
\begin{problem}
设函数 \(f(x)=\sin \omega x\)（\(\omega>0\)）. 已知 \(f(x_1)=-1\)，\(f(x_2)=1\)，且 \(x_1-x_2\) 的最小值为 \(\frac{\pi}{2}\)，则 \(\omega=\)（\quad）

\choices
    {1}
    {2}
    {3}
    {4}

\end{problem}

\begin{answer}
B
\end{answer}

\begin{solution}
根据题意，\(x_1\) 是 \(f(x)\) 的极小值点，\(x_2\) 是 \(f(x)\) 的极大值点，因此
\(x_1-x_2=\frac{T}{2}=\frac{\pi}{2}\).
于是 \(T=\pi\). 又
\(T=\frac{2\pi}{\omega}\)，
所以 \(\omega=2\). 故选 B.
\end{solution}
\begin{problem}
生物丰富度指数 \(d=\frac{S-1}{\ln N}\) 是河流水质的一个评价指标，其中 \(S,N\) 分别表示河流中的生物种类数与生物个体总数. 生物丰富度指数 \(d\) 越大，水质越好. 如果某河流治理前后的生物种类数 \(S\) 没有变化，生物个体总数由 \(N_1\) 变为 \(N_2\)，生物丰富度指数由 \(2.1\) 提高到 \(3.15\)，则（\quad）

\choices
    {\(3N_2=2N_1\)}
    {\(2N_2=3N_1\)}
    {\(N_2^2=N_1^3\)}
    {\(N_2^3=N_1^2\)}

\end{problem}

\begin{answer}
D
\end{answer}

\begin{solution}
根据题意，
\(\frac{S-1}{\ln N_1}=2.1\)，\(\frac{S-1}{\ln N_2}=3.15\).
消去 \(S-1\) 得
\(2.1\ln N_1=3.15\ln N_2\)，
即
\(2\ln N_1=3\ln N_2\).
因此
\(\ln N_1^2=\ln N_2^3\)，
所以 \(N_2^3=N_1^2\). 故选 D.
\end{solution}
\begin{problem}
如图，在四棱锥 \(P-ABCD\) 中，底面 \(ABCD\) 是边长为 4 的正方形，\(PA=PB=4\)，\(PC=PD=2\sqrt{2}\)，该棱锥的高为（\quad）

\bitmapfigure[max width=5.2cm,max totalheight=4.8cm]{img/2024/beijing/q8_fig1.png}

\choices
    {1}
    {2}
    {\(\sqrt{2}\)}
    {\(\sqrt{3}\)}

\end{problem}

\begin{answer}
D
\end{answer}

\begin{solution}
底面 \(ABCD\) 为正方形. 当相邻的棱长相等时，不妨设 \(PA=PB=AB=4\)，\(PC=PD=2\sqrt{2}\). 分别取 \(AB,CD\) 的中点 \(E,F\)，连接 \(PE,PF,EF\).

\bitmapfigure[max width=5.2cm,max totalheight=4.8cm]{img/2024/beijing/q8_solution_fig1.png}

由图中关系有 \(PE\perp AB\)，\(EF\perp AB\)，且 \(PE\cap EF=E\)，又 \(PE,EF\subset \text{平面 }PEF\)，所以 \(AB\perp \text{平面 }PEF\). 又 \(AB\subset \text{平面 }ABCD\)，所以 \(\text{平面 }PEF\perp \text{平面 }ABCD\). 过 \(P\) 作 \(EF\) 的垂线，垂足为 \(O\)，即 \(PO\perp EF\). 因为 \(\text{平面 }PEF\cap \text{平面 }ABCD=EF\)，且 \(PO\subset \text{平面 }PEF\)，所以 \(PO\perp \text{平面 }ABCD\)，即 \(PO\) 为棱锥的高.

由题意可得
\(PE=2\sqrt{3}\)，\(PF=2\)，\(EF=4\)，
于是
\(PE^2+PF^2=EF^2\)，
所以 \(PE\perp PF\)，从而
\(\frac{1}{2}PE\cdot PF=\frac{1}{2}PO\cdot EF\).
解得
\(PO=\frac{PE\cdot PF}{EF}=\sqrt{3}\).
当相对的棱长相等时，不妨设 \(PA=PC=4\)，\(PB=PD=2\sqrt{2}\). 因为
\(BD=4=\sqrt{PB^2+PD^2}\)，
此时不能形成三角形 \(PBD\)，这种情况不存在.

故选 D.
\end{solution}
\begin{problem}
已知 \((x_1,y_1)\)，\((x_2,y_2)\) 是函数 \(y=2^x\) 的图象上两个不同的点，则（\quad）

\choices
    {\(\log_2 \frac{y_1+y_2}{2}<\frac{x_1+x_2}{2}\)}
    {\(\log_2 \frac{y_1+y_2}{2}>\frac{x_1+x_2}{2}\)}
    {\(\log_2 \frac{y_1+y_2}{2}<x_1+x_2\)}
    {\(\log_2 \frac{y_1+y_2}{2}>x_1+x_2\)}

\end{problem}

\begin{answer}
B
\end{answer}

\begin{solution}
不妨设 \(x_1<x_2\). 因为函数 \(y=2^x\) 单调递增，所以 \(0<y_1<y_2\). 由均值不等式，
\(\frac{y_1+y_2}{2}>\sqrt{y_1y_2}=\sqrt{2^{x_1}2^{x_2}}=2^{\frac{x_1+x_2}{2}}\).
因为 \(\log_2 x\) 单调递增，所以
\(\log_2 \frac{y_1+y_2}{2}>\frac{x_1+x_2}{2}\)，
故 B 正确.

例如取 \(x_1=0\)，\(x_2=1\)，则 \(\log_2 \frac{y_1+y_2}{2}=\log_2 \frac{3}{2}<1=x_1+x_2\)，故 D 错误；取 \(x_1=-1\)，\(x_2=-2\)，则 \(\log_2 \frac{y_1+y_2}{2}=\log_2 \frac{3}{8}>-3=x_1+x_2\)，故 C 错误. 故选 B.
\end{solution}
\begin{problem}
已知 \(M=\left\{ (x,y)\mid y=x+t\left( x^2-x \right),\,1\le x\le 2,\,0\le t\le 1 \right\}\) 是平面直角坐标系中的点集. 设 \(d\) 是 \(M\) 中两点间距离的最大值，\(S\) 是 \(M\) 表示的图形的面积，则（\quad）

\choices
    {\(d=3\)，\(S<1\)}
    {\(d=3\)，\(S>1\)}
    {\(d=\sqrt{10}\)，\(S<1\)}
    {\(d=\sqrt{10}\)，\(S>1\)}

\end{problem}

\begin{answer}
C
\end{answer}

\begin{solution}
对任意给定的 \(x\in[1,2]\)，有 \(x^2-x=x(x-1)\ge 0\)，且 \(0\le t\le 1\)，所以
\(x\le x+t\left( x^2-x \right)\le x+\left( x^2-x \right)=x^2\).
因此 \(M\) 表示的图形是满足 \(y\le x^2\)，\(y\ge x\)，\(1\le x\le 2\) 的区域.

\bitmapfigure[max width=3.8cm,max totalheight=4.8cm]{img/2024/beijing/q10_solution_fig1.png}

如图，记 \(A(1,1)\)，\(B(2,2)\)，\(C(2,4)\). 则两点间距离最大值为
\(d=|AC|=\sqrt{(2-1)^2+(4-1)^2}=\sqrt{10}\).
又阴影面积小于三角形 \(ABC\) 的面积，
\(S<S_{\triangle ABC}=\frac{1}{2}\times 1\times 2=1\).
故选 C.
\end{solution}
\section{填空题}

\begin{problem}
抛物线 \(y^2=16x\) 的焦点坐标为 \fillinblank{}. 

\end{problem}

\begin{answer}
\((4,0)\)
\end{answer}

\begin{solution}
标准方程 \(y^2=2px\) 的抛物线焦点坐标为 \(\left( \frac{p}{2},0 \right)\). 由 \(2p=16\) 得 \(p=8\)，故焦点坐标为 \((4,0)\).
\end{solution}
\begin{problem}
在平面直角坐标系 \(xOy\) 中，角 \(\alpha\) 与角 \(\beta\) 均以 \(Ox\) 为始边，它们的终边关于原点对称. 若 \(\alpha\in \left[ \frac{\pi}{6},\frac{\pi}{3} \right]\)，则 \(\cos \beta\) 的最大值为 \fillinblank{}. 

\end{problem}

\begin{answer}
\(-\frac{1}{2}\)
\end{answer}

\begin{solution}
因为终边关于原点对称，所以
\(\beta=\alpha+\pi+2k\pi\)（\(k\in \Z\)）.
所以
\(\cos \beta=\cos \left( \alpha+\pi+2k\pi \right)=-\cos \alpha\).
当 \(\alpha\in \left[ \frac{\pi}{6},\frac{\pi}{3} \right]\) 时，
\(\cos \alpha\in \left[ \frac{1}{2},\frac{\sqrt{3}}{2} \right]\)，
故
\(\cos \beta\in \left[ -\frac{\sqrt{3}}{2},-\frac{1}{2} \right]\).
当且仅当 \(\alpha=\frac{\pi}{3}\) 时，\(\cos \beta\) 取得最大值 \(-\frac{1}{2}\).

因此最大值为 \(-\frac{1}{2}\).
\end{solution}
\begin{problem}
若直线 \(y=k(x-3)\) 与双曲线 \(\frac{x^2}{4}-y^2=1\) 只有一个公共点，则 \(k\) 的一个取值为 \fillinblank{}. 

\end{problem}

\begin{answer}
\(\frac{1}{2}\)（或 \(-\frac{1}{2}\)）
\end{answer}

\begin{solution}
联立 \(\frac{x^2}{4}-y^2=1\) 与 \(y=k(x-3)\)，化简得
\(\left( 1-4k^2 \right)x^2+24k^2x-36k^2-4=0\).
只有一个公共点，则需满足首项系数为 0 或判别式为 0. 由此解得
\(k=\pm \frac{1}{2}\).
故可填 \(\frac{1}{2}\) 或 \(-\frac{1}{2}\).
\end{solution}
\begin{problem}
汉代刘歆设计的“铜嘉量”是龠、合、升、斗、斛五量合一的标准量器，其中升量器、斗量器、斛量器的形状均可视为圆柱. 若升、斗、斛量器的容积成公比为 10 的等比数列，底面直径依次为 \(65\text{mm},325\text{mm},325\text{mm}\)，且斛量器的高为 \(230\text{mm}\)，则斗量器的高为 \fillinblank{}\(\text{mm}\)，升量器的高为 \fillinblank{}\(\text{mm}\). 

\end{problem}

\begin{answer}
23；\(\frac{115}{2}\)
\end{answer}

\begin{solution}
设升量器的高为 \(h_1\)，斗量器的高为 \(h_2\). 由容积成公比为 10 的等比数列可得
\(\frac{\pi\left( \frac{325}{2} \right)^2 h_2}{\pi\left( \frac{65}{2} \right)^2 h_1}=10\)，\(\frac{\pi\left( \frac{325}{2} \right)^2\times 230}{\pi\left( \frac{325}{2} \right)^2 h_2}=10\).
由第二式得
\(h_2=23\).
再代入第一式得
\(h_1=\frac{115}{2}\).
故答案为 23 和 \(\frac{115}{2}\).
\end{solution}
\begin{problem}
设 \(\{a_n\}\) 与 \(\{b_n\}\) 是两个不同的无穷数列，且都不是常数列. 记集合 \(M=\{k\mid a_k=b_k,\ k\in \N^*\}\)，给出下列 4 个结论：

\begin{circlelist}
\item 若 \(\{a_n\}\) 与 \(\{b_n\}\) 均为等差数列，则 \(M\) 中最多有 1 个元素；

\item 若 \(\{a_n\}\) 与 \(\{b_n\}\) 均为等比数列，则 \(M\) 中最多有 2 个元素；

\item 若 \(\{a_n\}\) 为等差数列，\(\{b_n\}\) 为等比数列，则 \(M\) 中最多有 3 个元素；

\item 若 \(\{a_n\}\) 为递增数列，\(\{b_n\}\) 为递减数列，则 \(M\) 中最多有 1 个元素.
\end{circlelist}

其中正确结论的序号是 \fillinblank{}. 

\end{problem}

\begin{answer}
\(\circled{1}\circled{3}\circled{4}\)
\end{answer}

\begin{solution}
\circled{1} 两个非常数等差数列的散点图都落在直线上，两条不同直线最多有一个交点，所以结论正确.

\circled{2} 取
\(a_n=2^{n-1}\)，\(b_n=-(-2)^{n-1}\)，
则当 \(n\) 为偶数时有 \(a_n=b_n\)，于是 \(M\) 中有无穷多个元素，故结论错误.

\circled{3} 设
\(a_n=kn+b\)（\(k\ne 0\)），\(b_n=Aq^n\)（\(Aq\ne 0,\ q\ne \pm 1\)）.
若 \(M\) 中至少有 4 个元素，则方程
\(Aq^n=kn+b\)
至少有 4 个不同的正整数解.

先看偶数解. 设 \(n=2m\)，则方程化为
\(A\left( q^2 \right)^m=2km+b\).
这个方程至多有两个偶数解；若恰有两个偶数解，则指数型一侧与线性一侧在该子问题中必须同向变化，从而有 \(Ak\ln q>0\). 否则二者反向变化，偶数解至多只有一个.

再看奇数解. 设 \(n=2m-1\)，则方程可化为
\(Aq\left( q^2 \right)^{m-1}=2km-k+b\).
用同样的方法可以得到：这个方程至多有两个奇数解；若恰有两个奇数解，则必须有 \(Ak\ln q<0\). 否则奇数解至多只有一个.

因此，不可能同时既有两个偶数解又有两个奇数解，所以方程 \(Aq^n=kn+b\) 不可能有 4 个不同的正整数解，即 \(M\) 中最多有 3 个元素，结论正确.

\circled{4} 若 \(\{a_n\}\) 递增、\(\{b_n\}\) 递减，则两列散点图分别呈上升趋势和下降趋势，最多只有一个交点，结论正确.

故答案为 \(\circled{1}\circled{3}\circled{4}\).
\end{solution}
\section{解答题}

\begin{problem}
在 \(\triangle ABC\) 中，内角 \(A,B,C\) 的对边分别为 \(a,b,c\)，\(\angle A\) 为钝角，\(a=7\)，\(\sin 2B=\frac{\sqrt{3}}{7}b\cos B\).

\begin{enumerate}
\item 求 \(\angle A\)；
\item 从条件 \(\circled{1}\)、条件 \(\circled{2}\)、条件 \(\circled{3}\) 这三个条件中选择一个作为已知，使得 \(\triangle ABC\) 存在，求 \(\triangle ABC\) 的面积.
\end{enumerate}

条件 \(\circled{1}\)：\(b=7\)；条件 \(\circled{2}\)：\(\cos B=\frac{13}{14}\)；条件 \(\circled{3}\)：\(c\sin A=\frac{5\sqrt{3}}{2}\).

\end{problem}

\begin{solution}
（1）由
\(\sin 2B=2\sin B\cos B=\frac{\sqrt{3}}{7}b\cos B\)
又 \(A\) 为钝角，所以 \(\cos B\ne 0\)，从而
\(2\sin B=\frac{\sqrt{3}}{7}b\).
由正弦定理，
\(\frac{b}{\sin B}=\frac{a}{\sin A}=\frac{7}{\sin A}\)，
于是
\(\sin A=\frac{\sqrt{3}}{2}\).
又因为 \(A\) 为钝角，所以
\(A=\frac{2\pi}{3}\).

（2）先判断三个条件.

若选条件 \(\circled{1}\)，则
\(\sin B=\frac{\sqrt{3}}{14}b=\frac{\sqrt{3}}{2}\).
因为 \(B\) 为三角形内角，所以 \(B=\frac{\pi}{3}\). 这时
\(A+B=\pi\)，
与三角形内角和矛盾，因此条件 \(\circled{1}\) 不可取.

若选条件 \(\circled{2}\)，则
\(\sin B=\sqrt{1-\left( \frac{13}{14} \right)^2}=\frac{3\sqrt{3}}{14}\).
由
\(2\sin B=\frac{\sqrt{3}}{7}b\)
得 \(b=3\). 又
\[
\sin C=\sin(A+B)=\sin \left( \frac{2\pi}{3}+B \right)=\frac{\sqrt{3}}{2}\cos B-\frac{1}{2}\sin B=\frac{5\sqrt{3}}{14}
\]
所以
\[
S_{\triangle ABC}=\frac{1}{2}ab\sin C=\frac{1}{2}\times 7\times 3\times \frac{5\sqrt{3}}{14}=\frac{15\sqrt{3}}{4}
\]

若选条件 \(\circled{3}\)，由
\(c\sin A=\frac{5\sqrt{3}}{2}\)
及 \(\sin A=\frac{\sqrt{3}}{2}\) 得 \(c=5\). 由正弦定理，
\(\frac{a}{\sin A}=\frac{c}{\sin C}\)，
即
\(\frac{7}{\frac{\sqrt{3}}{2}}=\frac{5}{\sin C}\)，
从而
\(\sin C=\frac{5\sqrt{3}}{14}\)，\(\cos C=\frac{11}{14}\).
于是
\[
\sin B=\sin(A+C)=\sin \left( \frac{2\pi}{3}+C \right)=\frac{\sqrt{3}}{2}\cos C-\frac{1}{2}\sin C=\frac{3\sqrt{3}}{14}
\]
故
\[
S_{\triangle ABC}=\frac{1}{2}ac\sin B=\frac{1}{2}\times 7\times 5\times \frac{3\sqrt{3}}{14}=\frac{15\sqrt{3}}{4}
\]
综上，选择条件 \(\circled{2}\) 或条件 \(\circled{3}\) 时，\(\triangle ABC\) 的面积均为 \(\frac{15\sqrt{3}}{4}\).
\end{solution}
\begin{problem}
如图，在四棱锥 \(P-ABCD\) 中，\(BC\parallel AD\)，\(AB=BC=1\)，\(AD=3\)，点 \(E\) 在 \(AD\) 上，且 \(PE\perp AD\)，\(PE=DE=2\).

\bitmapfigure[max width=5.1cm,max totalheight=4.8cm]{img/2024/beijing/q17_fig1.png}

\begin{enumerate}
\item 若 \(F\) 为线段 \(PE\) 中点，求证：\(BF\parallel \text{平面 }PCD\)；
\item 若 \(AB\perp \text{平面 }PAD\)，求平面 \(PAB\) 与平面 \(PCD\) 夹角的余弦值.
\end{enumerate}

\end{problem}

\begin{solution}
（1）取 \(PD\) 的中点为 \(S\)，连接 \(SF,SC\). 因为 \(F\) 为 \(PE\) 中点，所以
\(SF\parallel ED\)，\(SF=\frac{1}{2}ED=1\).
又 \(ED\parallel BC\)，且 \(ED=2BC\)，故
\(SF\parallel BC\)，\(SF=BC\).
于是四边形 \(SFBC\) 为平行四边形，从而
\(BF\parallel SC\).
因为 \(SC\subset \text{平面 }PCD\)，且 \(BF\not\subset \text{平面 }PCD\)，所以
\(BF\parallel \text{平面 }PCD\).

（2）由 \(ED=2\) 得 \(AE=1\). 又 \(AE\parallel BC\)，且 \(AE=BC\)，所以四边形 \(AECB\) 为平行四边形，从而
\(CE\parallel AB\).
因为 \(AB\perp \text{平面 }PAD\)，所以 \(CE\perp \text{平面 }PAD\). 又 \(PE,ED\subset \text{平面 }PAD\)，故
\(CE\perp PE\)，\(CE\perp ED\).
再结合 \(PE\perp ED\)，建立如图所示的空间直角坐标系.

\bitmapfigure[max width=5.4cm,max totalheight=4.8cm]{img/2024/beijing/q17_solution_fig1.png}

取
\(A(0,-1,0)\)，\(B(1,-1,0)\)，\(C(1,0,0)\)，\(D(0,2,0)\)，\(P(0,0,2)\).
则
\(\overrightarrow{PA}=(0,-1,-2)\)，\(\overrightarrow{PB}=(1,-1,-2)\)，\(\overrightarrow{PC}=(1,0,-2)\)，\(\overrightarrow{PD}=(0,2,-2)\).

设平面 \(PAB\) 的法向量为 \(\bs{m}=(x,y,z)\)，则由 \(\bs{m}\cdot \overrightarrow{PA}=0\) 且 \(\bs{m}\cdot \overrightarrow{PB}=0\)，可取 \(\bs{m}=(0,-2,1)\).

设平面 \(PCD\) 的法向量为 \(\bs{n}=(a,b,c)\)，则由 \(\bs{n}\cdot \overrightarrow{PC}=0\) 且 \(\bs{n}\cdot \overrightarrow{PD}=0\)，可取 \(\bs{n}=(2,1,1)\).

于是平面 \(PAB\) 与平面 \(PCD\) 夹角的余弦值为
\[
\frac{|\bs{m}\cdot \bs{n}|}{|\bs{m}||\bs{n}|}=\frac{|0\times 2+(-2)\times 1+1\times 1|}{\sqrt{5}\sqrt{6}}=\frac{1}{\sqrt{30}}=\frac{\sqrt{30}}{30}
\]
\end{solution}
\begin{problem}
某保险公司为了了解该公司某种保险产品的索赔情况，从合同险期限届满的保单中随机抽取 1000 份，记录并整理这些保单的索赔情况，获得数据如下表：

\begin{center}
\begin{tabular}{|c|c|c|c|c|c|}
\hline
赔偿次数 & 0 & 1 & 2 & 3 & 4 \\
\hline
单数 & 800 & 100 & 60 & 30 & 10 \\
\hline
\end{tabular}
\end{center}

假设：一份保单的保费为 \(0.4\) 万元；前 3 次索赔时，保险公司每次赔偿 \(0.8\) 万元；第四次索赔时，保险公司赔偿 \(0.6\) 万元. 假设不同保单的索赔次数相互独立. 用频率估计概率.

\begin{enumerate}
\item 估计一份保单索赔次数不少于 2 的概率；
\item 一份保单的毛利润定义为这份保单的保费与赔偿总金额之差.
    \begin{enumerate}
    \item 记 \(X\) 为一份保单的毛利润，估计 \(E(X)\)；
    \item 如果无索赔的保单的保费减少 \(4\%\)，有索赔的保单的保费增加 \(20\%\)，试比较这种情况下一份保单毛利润的数学期望估计值与（1）中 \(E(X)\) 估计值的大小（结论不要求证明）.
    \end{enumerate}
\end{enumerate}

\end{problem}

\begin{solution}
（1）设事件 \(A\) 表示“随机抽取一份保单，索赔次数不少于 2 次”，则
\(P(A)=\frac{60+30+10}{1000}=\frac{1}{10}\).

（2）（i）设赔付金额为 \(\xi\)，则 \(\xi\) 的可能取值为
\(0,\ 0.8,\ 1.6,\ 2.4,\ 3\).
由题设中的统计数据可得
\(P(\xi=0)=\frac{800}{1000}=\frac{4}{5}\)，\(P(\xi=0.8)=\frac{100}{1000}=\frac{1}{10}\)，
\(P(\xi=1.6)=\frac{60}{1000}=\frac{3}{50}\)，\(P(\xi=2.4)=\frac{30}{1000}=\frac{3}{100}\)，\(P(\xi=3)=\frac{10}{1000}=\frac{1}{100}\).
所以 \(E(\xi)=0\times \frac{4}{5}+0.8\times \frac{1}{10}+1.6\times \frac{3}{50}+2.4\times \frac{3}{100}+3\times \frac{1}{100}=0.278\).
因此
\(E(X)=0.4-0.278=0.122\)
（万元）.

（2）（ii）若保费调整后毛利润记为 \(Y\)，则新保费的数学期望估计值为
\(0.4\times \frac{4}{5}\times 96\%+0.4\times \frac{1}{5}\times 1.2=0.4032\).
所以
\(E(Y)=0.122+0.4032-0.4=0.1252\)
（万元）.

故
\(E(Y)>E(X)\).
即这种情况下一份保单毛利润的数学期望估计值大于（1）中 \(E(X)\) 的估计值.
\end{solution}
\begin{problem}
已知椭圆 \(E:\frac{x^2}{a^2}+\frac{y^2}{b^2}=1\)（\(a>b>0\)），以椭圆 \(E\) 的焦点和短轴端点为顶点的四边形是边长为 2 的正方形. 过点 \((0,t)\)（\(t>\sqrt{2}\)）且斜率存在的直线与椭圆 \(E\) 交于不同的两点 \(A,B\)，过点 \(A\) 和 \(C(0,1)\) 的直线 \(AC\) 与椭圆 \(E\) 的另一个交点为 \(D\).

\begin{enumerate}
\item 求椭圆 \(E\) 的方程及离心率；
\item 若直线 \(BD\) 的斜率为 0，求 \(t\) 的值.
\end{enumerate}

\end{problem}

\begin{solution}
（1）由题意，椭圆的短轴半长与焦距都等于 \(\sqrt{2}\)，即
\(b=c=\sqrt{2}\).
因此
\(a=\sqrt{b^2+c^2}=2\).
故椭圆方程为
\(\frac{x^2}{4}+\frac{y^2}{2}=1\)，
离心率为
\(e=\frac{c}{a}=\frac{\sqrt{2}}{2}\).

（2）直线 \(AB\) 的斜率不为 0，否则直线 \(AB\) 与椭圆无交点，矛盾. 设直线 \(AB\) 的方程为
\(y=kx+t\)（\(k\ne 0,\ t>\sqrt{2}\)），
设 \(A(x_1,y_1)\)，\(B(x_2,y_2)\). 联立 \(\frac{x^2}{4}+\frac{y^2}{2}=1\) 和 \(y=kx+t\)，化简并整理得
\(\left( 1+2k^2 \right)x^2+4ktx+2t^2-4=0\).
由题意，
\(\Delta=16k^2t^2-8\left( 2k^2+1 \right)\left( t^2-2 \right)=8\left( 4k^2+2-t^2 \right)>0\)，
即 \(k,t\) 应满足 \(4k^2+2-t^2>0\).

由韦达定理，
\(x_1+x_2=-\frac{4kt}{1+2k^2}\)，\(x_1x_2=\frac{2t^2-4}{1+2k^2}\).

\bitmapfigure[max width=4.5cm,max totalheight=4.8cm]{img/2024/beijing/q19_solution_fig1.png}

若直线 \(BD\) 的斜率为 0，由椭圆的对称性可设
\(D(-x_2,y_2)\).
直线 \(AD\) 的方程可写为
\(y=\frac{y_1-y_2}{x_1+x_2}(x-x_1)+y_1\).
因为 \(C(0,1)\) 在直线 \(AD\) 上，所以当 \(x=0\) 时有
\(1=\frac{x_1y_2+x_2y_1}{x_1+x_2}\).
代入 \(y_1=kx_1+t\)，\(y_2=kx_2+t\) 及韦达定理化简，得
\(\frac{4k(t^2-2)}{-4kt}+t=\frac{2}{t}=1\).
故
\(t=2\).
此时
\(4k^2+2-t^2=4k^2-2>0\)，
即
\(k<-\frac{\sqrt{2}}{2}\)或\(k>\frac{\sqrt{2}}{2}\).
因此 \(t=2\) 满足题意.
\end{solution}
\begin{problem}
设函数 \(f(x)=x+k\ln(1+x)\)（\(k\ne 0\)），直线 \(l\) 是曲线 \(y=f(x)\) 在点 \((t,f(t))\)（\(t>0\)）处的切线.

\begin{enumerate}
\item 当 \(k=-1\) 时，求 \(f(x)\) 的单调区间；
\item 求证：\(l\) 不经过点 \((0,0)\)；
\item 当 \(k=1\) 时，设点 \(A(t,f(t))\)（\(t>0\)），\(C(0,f(t))\)，\(O(0,0)\)，\(B\) 为 \(l\) 与 \(y\) 轴的交点，\(S_{\triangle ACO}\) 与 \(S_{\triangle ABO}\) 分别表示 \(\triangle ACO\) 与 \(\triangle ABO\) 的面积. 是否存在点 \(A\) 使得 \(2S_{\triangle ACO}=15S_{\triangle ABO}\) 成立？若存在，这样的点 \(A\) 有几个？
\end{enumerate}

（参考数据：\(1.09<\ln 3<1.10\)，\(1.60<\ln 5<1.61\)，\(1.94<\ln 7<1.95\)）

\end{problem}

\begin{solution}
（1）当 \(k=-1\) 时，
\(f(x)=x-\ln(1+x)\)，
所以
\(f'(x)=1-\frac{1}{1+x}=\frac{x}{1+x}\)（\(x>-1\)）.
当 \(x\in(-1,0)\) 时，\(f'(x)<0\)；当 \(x\in(0,+\infty)\) 时，\(f'(x)>0\). 因此 \(f(x)\) 在 \((-1,0)\) 上单调递减，在 \((0,+\infty)\) 上单调递增.

（2）由
\(f'(x)=1+\frac{k}{1+x}\)
知切线 \(l\) 在点 \((t,f(t))\) 处的斜率为
\(1+\frac{k}{1+t}\).
故切线方程为
\(y-f(t)=\left( 1+\frac{k}{1+t} \right)(x-t)\)（\(t>0\)）.
若 \(l\) 经过原点 \((0,0)\)，则
\(-f(t)=\left( 1+\frac{k}{1+t} \right)(-t)\)，
即
\(t+k\ln(1+t)=t+\frac{kt}{1+t}\).
于是
\(\ln(1+t)-\frac{t}{1+t}=0\).
设
\(F(t)=\ln(1+t)-\frac{t}{1+t}\).
则
\(F'(t)=\frac{1}{1+t}-\frac{1}{(1+t)^2}=\frac{t}{(1+t)^2}>0\)（\(t>0\)），
所以 \(F(t)\) 在 \((0,+\infty)\) 上单调递增，且
\(F(0)=0\).
因此 \(t>0\) 时恒有 \(F(t)>0\)，矛盾. 故 \(l\) 不经过点 \((0,0)\).

（3）当 \(k=1\) 时，
\(f(x)=x+\ln(1+x)\)，\(f'(x)=1+\frac{1}{1+x}=\frac{x+2}{x+1}>0\).
因为
\(S_{\triangle ACO}=\frac{1}{2}tf(t)\)，
且切线方程为
\(y-t-\ln(1+t)=\left( 1+\frac{1}{1+t} \right)(x-t)\)，
令 \(x=0\) 得 \(B\) 的纵坐标
\(q=\ln(1+t)-\frac{t}{1+t}\).
因此
\(S_{\triangle ABO}=\frac{1}{2}tq=\frac{1}{2}t\left( \ln(1+t)-\frac{t}{1+t} \right)\).

由条件 \(2S_{\triangle ACO}=15S_{\triangle ABO}\) 得
\(2tf(t)=15t\left( \ln(1+t)-\frac{t}{1+t} \right)\)，
化简为
\(13\ln(1+t)-2t-\frac{15t}{1+t}=0\).
设
\(h(t)=13\ln(1+t)-2t-\frac{15t}{1+t}\)（\(t>0\)）.
则
\(h'(t)=\frac{13}{1+t}-2-\frac{15}{(1+t)^2} =\frac{(-2t+1)(t-4)}{(1+t)^2}\).
所以 \(h(t)\) 在 \(\left( 0,\frac{1}{2} \right)\) 上单调递减，在 \(\left( \frac{1}{2},4 \right)\) 上单调递增，在 \((4,+\infty)\) 上单调递减.

又有
\(h\left( \frac{1}{2} \right)<0\)，
且
\(h(4)=13\ln 5-20>13\times 1.60-20=0.8>0\)，
\(h(24)=13\ln 25-48-\frac{15\times 24}{25} =26\ln 5-48-\frac{72}{5} <26\times 1.61-48-\frac{72}{5}<0\).
由零点存在定理并结合单调性，\(h(t)\) 在 \(\left( \frac{1}{2},4 \right)\) 上有一个零点，在 \((4,24)\) 上有一个零点.

\bitmapfigure[max width=3.6cm,max totalheight=4.8cm]{img/2024/beijing/q20_solution_fig1.png}

故满足条件的点 \(A\) 有 2 个.
\end{solution}
\begin{problem}
已知集合 \(M=\left\{ (i,j,k,w)\mid i\in\{1,2\},\ j\in\{3,4\},\ k\in\{5,6\},\ w\in\{7,8\},\ \text{且 }i+j+k+w\text{ 为偶数} \right\}\). 给定数列 \(A:a_1,a_2,\cdots,a_8\)，和序列 \(\Omega:T_1,T_2,\cdots,T_s\)，其中 \(T_t=(i_t,j_t,k_t,w_t)\in M\)（\(t=1,2,\cdots,s\)）. 对数列 \(A\) 进行如下变换：将 \(A\) 的第 \(i_1,j_1,k_1,w_1\) 项均加 1，其余项不变，得到的数列记作 \(T_1(A)\)；将 \(T_1(A)\) 的第 \(i_2,j_2,k_2,w_2\) 项均加 1，其余项不变，得到数列记作 \(T_2T_1(A)\)；……；以此类推，得到 \(T_s\cdots T_2T_1(A)\)，简记为 \(\Omega(A)\).

\begin{enumerate}
\item 给定数列 \(A:1,3,2,4,6,3,1,9\) 和序列 \(\Omega:(1,3,5,7)\)，\((2,4,6,8)\)，\((1,3,5,7)\)，写出 \(\Omega(A)\)；
\item 是否存在序列 \(\Omega\)，使得 \(\Omega(A)\) 为 \(a_1+2,a_2+6,a_3+4,a_4+2,a_5+8,a_6+2,a_7+4,a_8+4\)？若存在，写出一个符合条件的 \(\Omega\)；若不存在，请说明理由；
\item 若数列 \(A\) 的各项均为正整数，且 \(a_1+a_3+a_5+a_7\) 为偶数，求证：“存在序列 \(\Omega\)，使得 \(\Omega(A)\) 的各项都相等”的充要条件为“\(a_1+a_2=a_3+a_4=a_5+a_6=a_7+a_8\)”.
\end{enumerate}

\end{problem}

\begin{solution}
（1）按题意依次变换：
\(T_1(A):2,3,3,4,7,3,2,9\)，
\(T_2T_1(A):2,4,3,5,7,4,2,10\)，
\(T_3T_2T_1(A):3,4,4,5,8,4,3,10\).
故
\(\Omega(A):3,4,4,5,8,4,3,10\).

（2）解法一：假设存在符合条件的 \(\Omega\)，可知 \(\Omega(A)\) 的第 \(1,2\) 项之和为 \(a_1+a_2+s\)，第 \(3,4\) 项之和为 \(a_3+a_4+s\)，则 \((a_1+2)+(a_2+6)=a_1+a_2+s\) 且 \((a_3+4)+(a_4+2)=a_3+a_4+s\)，而该方程组无解，故假设不成立，所以不存在符合条件的 \(\Omega\).

解法二：根据题意可知，对任意序列，所得数列之和比原数列之和多 \(4\). 假设存在符合条件的 \(\Omega\)，且
\(\Omega(A):b_1,b_2,\cdots,b_8\).
因为
\(\frac{2+6+4+2+8+2+4+4}{4}=8\)，
所以序列 \(\Omega\) 共有 \(8\) 项. 根据题意，对 \(n=1,2,3,4\)，都有
\(\left( b_{2n-1}+b_{2n} \right)-\left( a_{2n-1}+a_{2n} \right)=8\).
检验可知，当 \(n=2,3\) 时，上式不成立，因此不存在符合条件的 \(\Omega\).

（3）解法一：我们设序列 \(T_s\cdots T_2T_1(A)\) 为 \(\{a_{s,n}\}\)（\(1\le n\le 8\)），特别规定 \(a_{0,n}=a_n\)（\(1\le n\le 8\)）.

必要性：

若存在序列 \(\Omega:T_1,T_2,\cdots,T_s\)，使得 \(\Omega(A)\) 的各项都相等，则
\[
a_{s,1}=a_{s,2}=a_{s,3}=a_{s,4}=a_{s,5}=a_{s,6}=a_{s,7}=a_{s,8}
\]
所以
\[
a_{s,1}+a_{s,2}=a_{s,3}+a_{s,4}=a_{s,5}+a_{s,6}=a_{s,7}+a_{s,8}
\]
根据 \(T_s\cdots T_2T_1(A)\) 的定义，显然有
\(a_{s,2j-1}+a_{s,2j}=a_{s-1,2j-1}+a_{s-1,2j}\)（\(j=1,2,3,4;\ s=1,2,\cdots\)）.
所以不断使用该式就得到
\[
a_1+a_2=a_3+a_4=a_5+a_6=a_7+a_8
\]
必要性得证.

充分性：

若
\[
a_1+a_2=a_3+a_4=a_5+a_6=a_7+a_8
\]
根据已知，\(a_1+a_3+a_5+a_7\) 为偶数，而
\(a_2+a_4+a_6+a_8=4\left( a_1+a_2 \right)-\left( a_1+a_3+a_5+a_7 \right)\)
也是偶数.

我们设 \(T_s\cdots T_2T_1(A)\) 是通过合法的序列 \(\Omega\) 的变换能得到的所有可能的数列 \(\Omega(A)\) 中，使得
\(\left|a_{s,1}-a_{s,2}\right|+\left|a_{s,3}-a_{s,4}\right|+\left|a_{s,5}-a_{s,6}\right|+\left|a_{s,7}-a_{s,8}\right|\)
最小的一个.

上面已经证明
\(a_{s,2j-1}+a_{s,2j}=a_{s-1,2j-1}+a_{s-1,2j}\)（\(j=1,2,3,4;\ s=1,2,\cdots\)），
从而由
\[
a_1+a_2=a_3+a_4=a_5+a_6=a_7+a_8
\]
可得
\[
a_{s,1}+a_{s,2}=a_{s,3}+a_{s,4}=a_{s,5}+a_{s,6}=a_{s,7}+a_{s,8}
\]
同时，由于 \(i_t+j_t+k_t+w_t\) 总是偶数，所以
\(a_{t,1}+a_{t,3}+a_{t,5}+a_{t,7}\)
和
\(a_{t,2}+a_{t,4}+a_{t,6}+a_{t,8}\)
的奇偶性保持不变，从而
\(a_{s,1}+a_{s,3}+a_{s,5}+a_{s,7}\)
和
\(a_{s,2}+a_{s,4}+a_{s,6}+a_{s,8}\)
都是偶数.

下面证明不存在 \(j=1,2,3,4\) 使得 \(\left|a_{s,2j-1}-a_{s,2j}\right|\ge 2\). 假设存在，根据对称性，不妨设
\(a_{s,1}-a_{s,2}\ge 2\).

情况 1：若
\(a_{s,3}-a_{s,4}+a_{s,5}-a_{s,6}+a_{s,7}-a_{s,8}=0\)，
则由
\(a_{s,1}+a_{s,3}+a_{s,5}+a_{s,7}\)
和
\(a_{s,2}+a_{s,4}+a_{s,6}+a_{s,8}\)
都是偶数可知 \(a_{s,1}-a_{s,2}\ge 4\). 对该数列连续作四次变换
\((2,3,5,8)\)，\((2,4,6,8)\)，\((2,3,6,7)\)，\((2,4,5,7)\)
后，新的
\(\left|a_{s+4,1}-a_{s+4,2}\right|+\left|a_{s+4,3}-a_{s+4,4}\right|+\left|a_{s+4,5}-a_{s+4,6}\right|+\left|a_{s+4,7}-a_{s+4,8}\right|\)
相比原来的
\(\left|a_{s,1}-a_{s,2}\right|+\left|a_{s,3}-a_{s,4}\right|+\left|a_{s,5}-a_{s,6}\right|+\left|a_{s,7}-a_{s,8}\right|\)
减少 \(4\)，这与最小性矛盾.

情况 2：若
\(a_{s,3}-a_{s,4}+a_{s,5}-a_{s,6}+a_{s,7}-a_{s,8}>0\)，
不妨设 \(a_{s,3}-a_{s,4}>0\).

情况 2-1：如果 \(a_{s,3}-a_{s,4}\ge 1\)，则对该数列连续作两次变换
\((2,4,5,7)\)，\((2,4,6,8)\)
后，新的
\(\left|a_{s+2,1}-a_{s+2,2}\right|+\left|a_{s+2,3}-a_{s+2,4}\right|+\left|a_{s+2,5}-a_{s+2,6}\right|+\left|a_{s+2,7}-a_{s+2,8}\right|\)
相比原来的
\(\left|a_{s,1}-a_{s,2}\right|+\left|a_{s,3}-a_{s,4}\right|+\left|a_{s,5}-a_{s,6}\right|+\left|a_{s,7}-a_{s,8}\right|\)
至少减少 \(2\)，这与最小性矛盾.

情况 2-2：如果 \(a_{s,4}-a_{s,3}\ge 1\)，则对该数列连续作两次变换
\((2,3,5,8)\)，\((2,3,6,7)\)
后，新的
\(\left|a_{s+2,1}-a_{s+2,2}\right|+\left|a_{s+2,3}-a_{s+2,4}\right|+\left|a_{s+2,5}-a_{s+2,6}\right|+\left|a_{s+2,7}-a_{s+2,8}\right|\)
相比原来的
\(\left|a_{s,1}-a_{s,2}\right|+\left|a_{s,3}-a_{s,4}\right|+\left|a_{s,5}-a_{s,6}\right|+\left|a_{s,7}-a_{s,8}\right|\)
至少减少 \(2\)，这与最小性矛盾.

这就说明无论如何都会导致矛盾，所以对任意的 \(j=1,2,3,4\) 都有
\(\left|a_{s,2j-1}-a_{s,2j}\right|\le 1\).

假设存在 \(j=1,2,3,4\) 使得 \(\left|a_{s,2j-1}-a_{s,2j}\right|=1\)，则
\(a_{s,2j-1}+a_{s,2j}\)
是奇数，所以
\[
a_{s,1}+a_{s,2}=a_{s,3}+a_{s,4}=a_{s,5}+a_{s,6}=a_{s,7}+a_{s,8}
\]
都是奇数，设为 \(2N+1\). 则此时对任意 \(j=1,2,3,4\)，由 \(\left|a_{s,2j-1}-a_{s,2j}\right|\le 1\) 可知必有
\(\{a_{s,2j-1},a_{s,2j}\}=\{N,N+1\}\).
而
\(a_{s,1}+a_{s,3}+a_{s,5}+a_{s,7}\)
和
\(a_{s,2}+a_{s,4}+a_{s,6}+a_{s,8}\)
都是偶数，故集合
\(\{m\mid a_{s,m}=N\}\)
中的四个元素 \(i,j,k,w\) 之和为偶数. 对该数列进行一次变换 \((i,j,k,w)\)，则该数列成为常数列，新的
\(\left|a_{s+1,1}-a_{s+1,2}\right|+\left|a_{s+1,3}-a_{s+1,4}\right|+\left|a_{s+1,5}-a_{s+1,6}\right|+\left|a_{s+1,7}-a_{s+1,8}\right|\)
等于零，比原来的
\(\left|a_{s,1}-a_{s,2}\right|+\left|a_{s,3}-a_{s,4}\right|+\left|a_{s,5}-a_{s,6}\right|+\left|a_{s,7}-a_{s,8}\right|\)
更小，这与最小性矛盾.

综上，只可能
\(a_{s,2j-1}=a_{s,2j}\)（\(j=1,2,3,4\)），
而
\[
a_{s,1}+a_{s,2}=a_{s,3}+a_{s,4}=a_{s,5}+a_{s,6}=a_{s,7}+a_{s,8}
\]
故 \(\{a_{s,n}\}=\Omega(A)\) 是常数列，充分性得证.

解法二：由题意可知，\(\Omega\) 中序列的顺序不影响 \(\Omega(A)\) 的结果，且 \((a_1,a_2)\)、\((a_3,a_4)\)、\((a_5,a_6)\)、\((a_7,a_8)\) 相对于序列也是无序的.

（1）若
\[
a_1+a_2=a_3+a_4=a_5+a_6=a_7+a_8
\]
不妨设
\(a_1\le a_3\le a_5\le a_7\)，
则
\(a_2\ge a_4\ge a_6\ge a_8\).

\circled{1} 当 \(a_1=a_3=a_5=a_7\) 时，则 \(a_8=a_6=a_4=a_2\). 分别执行 \(a_1\) 个序列 \((2,4,6,8)\)、\(a_2\) 个序列 \((1,3,5,7)\)，可得
\(a_1+a_2,\ a_1+a_2,\ a_1+a_2,\ a_1+a_2,\ a_1+a_2,\ a_1+a_2,\ a_1+a_2,\ a_1+a_2\)，
为常数列.

\circled{2} 当 \(a_1,a_3,a_5,a_7\) 中有且仅有三个数相等时，不妨设 \(a_1=a_3=a_5\)，则 \(a_2=a_4=a_6\)，即
\(a_1,a_2,a_1,a_2,a_1,a_2,a_7,a_8\).
分别执行 \(a_2\) 个序列 \((1,3,5,7)\)、\(a_7\) 个序列 \((2,4,6,8)\)，可得
\(a_1+a_2,\ a_2+a_7,\ a_1+a_2,\ a_2+a_7,\ a_1+a_2,\ a_2+a_7,\ a_2+a_7,\ a_7+a_8\).
即
\(a_1+a_2,\ a_2+a_7,\ a_1+a_2,\ a_2+a_7,\ a_1+a_2,\ a_2+a_7,\ a_2+a_7,\ a_1+a_2\).
因为 \(a_1+a_3+a_5+a_7\) 为偶数，即 \(3a_1+a_7\) 为偶数，所以 \(a_1,a_7\) 的奇偶性相同，从而
\(\frac{a_7-a_1}{2}\in \N^*\).
分别执行 \(\frac{a_7-a_1}{2}\) 个序列 \((1,3,5,7)\)、\((1,3,6,8)\)、\((2,3,5,8)\)、\((1,4,5,8)\)，可得
\(\frac{3a_7+2a_2-a_1}{2}\)，\(\frac{3a_7+2a_2-a_1}{2}\)，\(\frac{3a_7+2a_2-a_1}{2}\)，\(\frac{3a_7+2a_2-a_1}{2}\)，
\(\frac{3a_7+2a_2-a_1}{2}\)，\(\frac{3a_7+2a_2-a_1}{2}\)，\(\frac{3a_7+2a_2-a_1}{2}\)，\(\frac{3a_7+2a_2-a_1}{2}\)，
为常数列.

\circled{3} 若 \(a_1=a_3<a_5=a_7\)，则 \(a_2=a_4>a_6=a_8\)，即
\(a_1,a_2,a_1,a_2,a_5,a_6,a_5,a_6\).
分别执行 \(a_5\) 个序列 \((1,3,6,8)\)、\(a_1\) 个序列 \((2,4,5,7)\)，可得
\(a_1+a_5,\ a_1+a_2,\ a_1+a_5,\ a_1+a_2,\ a_1+a_5,\ a_5+a_6,\ a_1+a_5,\ a_5+a_6\).
因为 \(a_1+a_2=a_5+a_6\)，可得
\(a_1+a_5,\ a_1+a_2,\ a_1+a_5,\ a_1+a_2,\ a_1+a_5,\ a_1+a_2,\ a_1+a_5,\ a_1+a_2\)，
即转为 \(\circled{1}\).

\circled{4} 当 \(a_1,a_3,a_5,a_7\) 中有且仅有两个数相等时，不妨设 \(a_1=a_3\)，则 \(a_2=a_4\)，即
\(a_1,a_2,a_1,a_2,a_5,a_6,a_7,a_8\).
分别执行 \(a_1\) 个 \((2,4,5,7)\)、\(a_5\) 个 \((1,3,6,8)\)，可得
\(a_1+a_5,\ a_1+a_2,\ a_1+a_5,\ a_1+a_2,\ a_1+a_5,\ a_5+a_6,\ a_1+a_7,\ a_5+a_8\).
且 \(a_1+a_2=a_5+a_6\)，可得
\(a_1+a_5,\ a_1+a_2,\ a_1+a_5,\ a_1+a_2,\ a_1+a_5,\ a_1+a_2,\ a_1+a_7,\ a_5+a_8\)，
即转为 \(\circled{2}\).

\circled{5} 若 \(a_1<a_3<a_5<a_7\)，则 \(a_2>a_4>a_6>a_8\)，即
\(a_1,a_2,a_3,a_4,a_5,a_6,a_7,a_8\).
分别执行 \(a_1\) 个 \((2,3,5,8)\)、\(a_3\) 个 \((1,4,6,7)\)，可得
\(a_1+a_3,\ a_1+a_2,\ a_1+a_3,\ a_3+a_4,\ a_1+a_5,\ a_3+a_6,\ a_3+a_7,\ a_1+a_8\).
且 \(a_1+a_2=a_3+a_4\)，可得
\(a_1+a_3,\ a_1+a_2,\ a_1+a_3,\ a_1+a_2,\ a_1+a_5,\ a_3+a_6,\ a_3+a_7,\ a_1+a_8\)，
即转为 \(\circled{3}\).

由上所述，若
\(a_1+a_2=a_3+a_4=a_5+a_6=a_7+a_8\)，
则存在序列 \(\Omega\)，使得 \(\Omega(A)\) 为常数列.

（2）若存在序列 \(\Omega\)，使得 \(\Omega(A)\) 为常数列，因为对任意
\(\Omega(A):b_1,b_2,\cdots,b_8\)，
均有
\[
\left( b_1+b_2 \right)-\left( a_1+a_2 \right)=\left( b_3+b_4 \right)-\left( a_3+a_4 \right)=\left( b_5+b_6 \right)-\left( a_5+a_6 \right)=\left( b_7+b_8 \right)-\left( a_7+a_8 \right)
\]
若 \(\Omega(A)\) 为常数列，则
\[
b_1+b_2=b_3+b_4=b_5+b_6=b_7+b_8
\]
所以
\[
a_1+a_2=a_3+a_4=a_5+a_6=a_7+a_8
\]

综上，“存在序列 \(\Omega\)，使得 \(\Omega(A)\) 的各项都相等”的充要条件为
\[
a_1+a_2=a_3+a_4=a_5+a_6=a_7+a_8
\]
\end{solution}
